% =====================================================================
%  PAQUETES Y FORMATO GENERAL
% =====================================================================

% =====================================================================
%  Formato según la "Guía para la elaboración de trabajos de grado"
%  (Carreras de Ing. Informática e Ing. de Sistemas, FCyT-UMSS, 2013):
%   - Fuente Arial 11 (párrafos), 12 MAYÚSCULAS (títulos), 11 (subtítulos)
%   - Interlineado sencillo, espaciado posterior 6 pt, justificado, sin sangría
%   - Carta; márgenes izq. 3 cm, der. 2 cm, sup. 2 cm, inf. 2 cm
%   - Número de página en el pie, centrado
%   - Figuras y cuadros: número, descripción y fuente DEBAJO; centrados y con recuadro
% =====================================================================

% ---------- Fuente Arial ----------
% Con XeLaTeX/LuaLaTeX se usa Arial real (o TeX Gyre Heros si no está instalada).
% Con pdfLaTeX se usa Helvetica, de métricas equivalentes a Arial.
\usepackage{iftex}
\ifPDFTeX
  \usepackage[utf8]{inputenc}
  \usepackage[T1]{fontenc}
  \usepackage[scaled=1]{helvet}
  \renewcommand{\familydefault}{\sfdefault}
\else
  \usepackage{fontspec}
  \IfFontExistsTF{Arial}{\setmainfont{Arial}\setsansfont{Arial}}
                       {\setmainfont{TeX Gyre Heros}\setsansfont{TeX Gyre Heros}}
\fi
\usepackage[spanish, es-tabla, es-noshorthands, es-nodecimaldot]{babel}
\usepackage{csquotes}

% ---------- Página carta y márgenes ----------
\usepackage[letterpaper, top=2cm, bottom=2cm, left=3cm, right=2cm,
            headheight=14pt, footskip=1cm]{geometry}

% ---------- Gráficos, colores y diagramas ----------
\usepackage{graphicx}
\usepackage[table]{xcolor}
\usepackage{tikz}
\usetikzlibrary{positioning, arrows.meta, shapes.geometric, shapes.misc,
                fit, backgrounds, calc, shadows, shapes.multipart}
\usepackage{pgfgantt}

% ---------- Tablas ----------
\usepackage{array}
\usepackage{booktabs}
\usepackage{tabularx}
\usepackage{longtable}
\usepackage{needspace}
\setlength{\LTpost}{2pt}
\usepackage{multirow}
\usepackage{makecell}
\renewcommand{\arraystretch}{1.25}
% Tablas con cuadrícula: bordes en todas las celdas y encabezado sombreado
\definecolor{encabezado}{gray}{0.88}
\setlength{\arrayrulewidth}{0.5pt}
\newcolumntype{L}[1]{>{\raggedright\arraybackslash}p{#1}}
\newcolumntype{C}[1]{>{\centering\arraybackslash}p{#1}}
\newcolumntype{Y}{>{\raggedright\arraybackslash}X}

% ---------- Listas y leyendas ----------
\usepackage{enumitem}
\setlist{itemsep=2pt, topsep=4pt, parsep=0pt}
% Leyendas DEBAJO de figuras y cuadros: "Figura 2-1. Descripción. Fuente: ..."
\usepackage[font=small, labelfont=bf, labelsep=period, justification=centering,
            figureposition=bottom, tableposition=bottom, skip=6pt]{caption}
% Recuadro de las figuras
\setlength{\fboxsep}{0pt}
\setlength{\fboxrule}{0.5pt}
\usepackage{float}
\usepackage{amsmath, amssymb}

% ---------- Código fuente ----------
\usepackage{listings}
\definecolor{codegreen}{rgb}{0,0.5,0.2}
\definecolor{codegray}{rgb}{0.5,0.5,0.5}
\definecolor{codepurple}{rgb}{0.58,0,0.82}
\definecolor{backcolour}{rgb}{0.96,0.97,0.96}

\lstdefinelanguage{TypeScript}{
  keywords={abstract, any, as, async, await, boolean, break, case, catch,
    class, const, continue, default, else, enum, export, extends, false,
    finally, for, from, function, if, implements, import, in, interface,
    let, new, null, number, of, private, public, readonly, return, string,
    switch, this, throw, true, try, type, typeof, undefined, void, while},
  sensitive=true,
  comment=[l]{//},
  morecomment=[s]{/*}{*/},
  morestring=[b]",
  morestring=[b]',
  morestring=[b]`
}

\lstdefinelanguage{FirestoreRules}{
  keywords={rules_version, service, match, allow, if, function, return,
    let, read, write, create, update, delete, true, false, is, in},
  sensitive=true,
  comment=[l]{//},
  morestring=[b]'
}

\lstdefinestyle{tesis}{
  backgroundcolor=\color{backcolour},
  commentstyle=\color{codegreen}\itshape,
  keywordstyle=\color{magenta!70!black}\bfseries,
  numberstyle=\tiny\color{codegray},
  stringstyle=\color{codepurple},
  basicstyle=\ttfamily\scriptsize,
  breakatwhitespace=false,
  breaklines=true,
  captionpos=b,
  keepspaces=true,
  numbers=left,
  numbersep=6pt,
  showspaces=false,
  showstringspaces=false,
  showtabs=false,
  tabsize=2,
  frame=single,
  rulecolor=\color{black!20},
  xleftmargin=1.5em,
  framexleftmargin=1.2em,
  extendedchars=true,
  literate=
    {á}{{\'a}}1 {é}{{\'e}}1 {í}{{\'i}}1 {ó}{{\'o}}1 {ú}{{\'u}}1
    {Á}{{\'A}}1 {É}{{\'E}}1 {Í}{{\'I}}1 {Ó}{{\'O}}1 {Ú}{{\'U}}1
    {ñ}{{\~n}}1 {Ñ}{{\~N}}1 {ü}{{\"u}}1 {¿}{{?`}}1 {¡}{{!`}}1
    {→}{{$\rightarrow$}}1 {·}{{$\cdot$}}1 {…}{{...}}3
    {—}{{--}}2 {─}{{-}}1 {│}{{|}}1 {├}{{+}}1 {└}{{+}}1 {┌}{{+}}1
    {┐}{{+}}1 {┘}{{+}}1 {┬}{{+}}1 {┴}{{+}}1 {┼}{{+}}1 {┤}{{+}}1
    {̀}{{\textbackslash{}u0300}}6 {ͯ}{{\textbackslash{}u036f}}6
}
\lstset{style=tesis}
\renewcommand{\lstlistingname}{Código}
\renewcommand{\lstlistlistingname}{ÍNDICE DE CÓDIGOS FUENTE}

% ---------- Hipervínculos ----------
\usepackage[hyphens]{url}
\usepackage{hyperref}
\hypersetup{
  colorlinks=true,
  linkcolor=black,
  urlcolor=black,
  citecolor=black,
  filecolor=black,
  pdftitle={Sistema inteligente de entrenamiento en gimnasios},
  pdfauthor={Erika Chino Alanoca}
}

% ---------- Bibliografía (estilo autor-año, compatible con APA) ----------
\usepackage[round, authoryear]{natbib}
\setcitestyle{aysep={,}}

% ---------- Interlineado y párrafos ----------
\usepackage{setspace}
\singlespacing                   % interlineado sencillo
\setlength{\parindent}{0pt}      % sin sangría
\setlength{\parskip}{6pt}        % espaciado posterior 6 pt

% ---------- Encabezados y pies de página ----------
\usepackage{fancyhdr}
\pagestyle{fancy}
\fancyhf{}
\fancyfoot[C]{\thepage}         % número de página centrado en el pie
\renewcommand{\headrulewidth}{0pt}
\fancypagestyle{plain}{\fancyhf{}\fancyfoot[C]{\thepage}\renewcommand{\headrulewidth}{0pt}}

% ---------- Numeración: capítulos en romanos, secciones en arábigos ----------
\renewcommand{\thechapter}{\Roman{chapter}}
\renewcommand{\thesection}{\arabic{chapter}.\arabic{section}}
\renewcommand{\thefigure}{\arabic{chapter}-\arabic{figure}}
\renewcommand{\thetable}{\arabic{chapter}-\arabic{table}}
\AtBeginDocument{\renewcommand{\thelstlisting}{\arabic{chapter}-\arabic{lstlisting}}}
\renewcommand{\theequation}{\arabic{chapter}.\arabic{equation}}
\setcounter{secnumdepth}{3}
\setcounter{tocdepth}{2}

% ---------- Títulos y subtítulos ----------
% Títulos (capítulos): 12 pt, negrilla, MAYÚSCULAS, centrados, 6 pt antes y después.
% Subtítulos: 11 pt, negrilla, alineados a la izquierda, 6 pt antes y después.
\usepackage{titlesec}
\newcommand{\fuentetitulo}{\fontsize{12}{14}\selectfont}
\titleformat{\chapter}[block]
  {\normalfont\fuentetitulo\bfseries\centering}
  {\MakeUppercase{\chaptertitlename}\ \thechapter:\ }
  {0pt}
  {\MakeUppercase}
\titlespacing*{\chapter}{0pt}{0pt}{6pt}
\titleformat{\section}{\normalfont\normalsize\bfseries}{\thesection.}{0.5em}{}
\titleformat{\subsection}{\normalfont\normalsize\bfseries}{\thesubsection.}{0.5em}{}
\titleformat{\subsubsection}{\normalfont\normalsize\bfseries}{\thesubsubsection.}{0.5em}{}
\titlespacing*{\section}{0pt}{6pt}{0pt}
\titlespacing*{\subsection}{0pt}{6pt}{0pt}
\titlespacing*{\subsubsection}{0pt}{6pt}{0pt}
\titlespacing*{\paragraph}{0pt}{6pt}{0.5em}

% ---------- Nombres en español ----------
\addto\captionsspanish{%
  \renewcommand{\contentsname}{ÍNDICE GENERAL}%
  \renewcommand{\listfigurename}{ÍNDICE DE FIGURAS}%
  \renewcommand{\listtablename}{ÍNDICE DE TABLAS}%
  \renewcommand{\appendixname}{Anexo}%
  \renewcommand{\chaptername}{Capítulo}%
}
% En los anexos la etiqueta es "ANEXO A -- ..."
\usepackage{etoolbox}
\appto\appendix{%
  \renewcommand{\thechapter}{\Alph{chapter}}%
  \renewcommand{\thesection}{\Alph{chapter}.\arabic{section}}%
  \renewcommand{\thefigure}{\Alph{chapter}-\arabic{figure}}%
  \renewcommand{\thetable}{\Alph{chapter}-\arabic{table}}%
  \renewcommand{\thelstlisting}{\Alph{chapter}.\arabic{lstlisting}}%
}

% ---------- Ancho de números en los índices ----------
\usepackage{tocloft}
\renewcommand{\cftchapfont}{\bfseries}
\renewcommand{\cftchappagefont}{\bfseries}
\setlength{\cftchapnumwidth}{2.6em}
\setlength{\cftsecnumwidth}{2.6em}
\setlength{\cftsubsecnumwidth}{3.4em}
\setlength{\cftfignumwidth}{2.8em}
\setlength{\cfttabnumwidth}{2.8em}
% Títulos de los índices: 12 pt, negrilla, centrados, sin salto de página propio
\renewcommand{\cfttoctitlefont}{\hfill\fuentetitulo\bfseries}
\renewcommand{\cftaftertoctitle}{\hfill}
\renewcommand{\cftloftitlefont}{\hfill\fuentetitulo\bfseries}
\renewcommand{\cftafterloftitle}{\hfill}
\renewcommand{\cftlottitlefont}{\hfill\fuentetitulo\bfseries}
\renewcommand{\cftafterlottitle}{\hfill}
\setlength{\cftbeforetoctitleskip}{0pt}
\setlength{\cftaftertoctitleskip}{6pt}
\setlength{\cftbeforeloftitleskip}{12pt}
\setlength{\cftafterloftitleskip}{6pt}
\setlength{\cftbeforelottitleskip}{12pt}
\setlength{\cftafterlottitleskip}{6pt}
\setlength{\cftbeforechapskip}{4pt}
\setlength{\emergencystretch}{2em}
